% Generated draft: requires mathematical and source review.
\documentclass[11pt]{article}
\usepackage[margin=25mm]{geometry}
\usepackage{fontspec}
\setmainfont{Latin Modern Roman}
\usepackage{amsmath,amssymb,mathtools,unicode-math}
\setmathfont{Latin Modern Math}
\usepackage{enumitem,needspace,xurl,hyperref}
\setlength{\parindent}{0pt}
\setlength{\parskip}{5pt}
\setlength{\emergencystretch}{3em}
\newcommand{\N}{\mathbb{N}}
\newcommand{\Z}{\mathbb{Z}}
\newcommand{\Q}{\mathbb{Q}}
\newcommand{\R}{\mathbb{R}}
\newcommand{\C}{\mathbb{C}}
\newcommand{\F}{\mathbb{F}}
\newcommand{\A}{\mathbb{A}}
\newcommand{\PP}{\mathbb{P}}
\newcommand{\E}{\mathbb{E}}
\newcommand{\eps}{\varepsilon}
\newcommand{\dd}{\,\mathrm{d}}
\begin{document}
\section*{UnsolvedMath ID: 30001551}
\textbf{Source problem code:} catalogue entry
\subsection{Definitions and mathematical statement}
Let \(X=\{x_1,x_2,x_3\}\) be three formal unknowns. Under the constant-free interpretation of word equations, an equation is a pair \(U=V\) with \(U,V\in X^*\), where \(X^*\) is the free monoid with identity the empty word \(\varepsilon\). For a target alphabet \(A\), a substitution is a monoid homomorphism \(h:X^*\to A^*\); thus empty images are permitted under this convention. For a system \(S\), define
\[
\operatorname{Sol}_A(S)=\{h:X^*\to A^*\text{ a monoid homomorphism}:h(U)=h(V)\text{ for every }(U=V)\in S\}.
\]
A three-equation system \(S=\{E_1,E_2,E_3\}\) is independent if no proper subsystem has the same solution set, equivalently if
\[
\operatorname{Sol}_A(S)\subsetneq\operatorname{Sol}_A(S\setminus\{E_i\})\qquad(i=1,2,3).
\]
In particular, for each \(i\) there must be a substitution satisfying the other two equations and failing \(E_i\); these witnesses need not be the same substitution.

Under the usual common-power convention, a substitution \(h\) is periodic if there exist \(p\in A^*\) and integers \(n_1,n_2,n_3\geq0\) such that
\[
h(x_j)=p^{n_j}\qquad(j=1,2,3).
\]
It is nonperiodic if no such common word and exponents exist.

The supplied question asks whether there exists an independent system of exactly three equations in these three unknowns that admits a nonperiodic solution. With the conventions above and a specified target alphabet \(A\), this is the existence assertion
\[
\exists S\subseteq X^*,\!\times X^*:\quad |S|=3,\quad
\bigwedge_{i=1}^{3}\operatorname{Sol}_A(S)\subsetneq\operatorname{Sol}_A(S\setminus\{E_i\}),\quad
\exists h\in\operatorname{Sol}_A(S)\ \forall p\in A^*\ \forall(n_1,n_2,n_3)\in\Z_{\geq0}^{3}\ \exists j\in\{1,2,3\}:h(x_j)\ne p^{n_j},
\]
where \(S=\{E_1,E_2,E_3\}\). Here the pair notation for equations means \(S\subseteq X^*\times X^*\).

\textit{Unresolved qualifications.} The source does not specify the target alphabet or its quantification, whether constants are allowed in equations, whether substitutions must be nonerasing, or an explicit definition of nonperiodicity. The displayed formulation uses constant-free equations, possibly erasing substitutions, and the common-power definition; it is conditional on those being the intended conventions. If constants are intended, their alphabet and the requirement that substitutions fix them must be specified.
\subsection{Short English statement}
Does an independent system of three word equations in three unknowns have a solution whose three assigned words are not all powers of one common word? Independence means that removing any equation enlarges the solution set. The source leaves alphabet, constant, and substitution conventions implicit, so the formalization requires review.
\subsection{Sources}
\begin{itemize}
\item Supplied local problem registry: permanent UnsolvedMath ID \(30001551\), source problem code \texttt{catalogue entry}, title \emph{Independent Word Equations with Nonperiodic Solutions}; collection: Oberwolfach Reports Open Problems.
\item Supplied source citation: \emph{Mini-Workshop: Combinatorics on Words (2011)}, Oberwolfach Reports, pp.\ 2195--2244, DOI \texttt{10.4171/owr/2010/37}. Source URL and supplied literature source labelled ``original workshop report'': \url{https://doi.org/10.4171/owr/2010/37}.
\item The background attributes the record to workshop open-problem discussions documented in Oberwolfach Reports. No individual proposer is supplied. The proposed year and year tag are 2010, whereas the supplied report title says 2011; this discrepancy is unresolved.
\item Literature assessment dated exactly \texttt{2026-08-21}: the supplied curation labels the problem open and states that the independent three-equation nonperiodic-solution problem remains unresolved, with later work providing bounds but no definitive answer. This is supplied dated curation, not independently verified. The only listed literature source is the report linked above; no references for the asserted later bounds are supplied.
\item The supplied record marks the statement as ``exact'' and explains that it was internally readable and retained byte-for-byte. This is a supplied provenance label, not independent verification of the report or of mathematical correctness.
\end{itemize}
\end{document}
